\section{Objective and Theorem Ladder}
\label{sec:v8-objective}
\label{sec:v8-theorem-ladder}

\subsection{The C-TreePO Objective}

Section~\ref{sec:v8-ctree-math} gives the structural goal: learn one state map
\(g\) whose node states can replace the manifesto spans they represent for
oracle-compatible readouts. C-TreePO turns that structural goal into an
optimization problem for the same learned \(g\). The root should be useful for
the document-level expert target, and the internal states should behave like
valid manifesto policy states.

Let \(L_{\mathrm{root}}(g)\) be the full-document root loss, measured against
the available expert target. For the local channel, attach one node loss to
each \(v\in V(T)\). This node loss collects the checks that apply at \(v\):
C1 for a leaf, C2 for the produced state, and C3 for an internal merge. If
every check could be measured by the target oracle, the quantity at node \(v\)
would be
\[
  \ell_v^\ast
  =
  \ell_{\mathrm{law}}(v;f^\ast,g).
\]
The star marks the target quantity: the local-law loss as it would be measured
by the expert oracle, not by the cheaper proxy used during dense training.

Let \(d(v)\) be the depth of node \(v\), with the root at depth \(0\), and let
\(\gamma\ge 0\) be a declared depth-discount weight. The oracle local-law
channel is
\begin{equation}
\label{eq:v8-local-law-estimand}
  L_{\mathrm{law}}^\ast(g,T)
  =
  \sum_{v\in V(T)}\gamma^{d(v)}\ell_v^\ast .
\end{equation}
The theorem-facing objective is
\begin{equation}
\label{eq:v8-main-objective}
  J_\lambda^\ast(g)
  =
  (1-\lambda)L_{\mathrm{root}}(g)
  +
  \lambda L_{\mathrm{law}}^\ast(g,T),
  \qquad 0\le\lambda\le 1 .
\end{equation}
Here \(\lambda\) is the analyst-chosen local-law share. The root term names
the manifesto score we ultimately report. The local term names the
state-validity pressure that makes the theorem relevant to the fitted tree.

\subsection{Three Claim Stacks}

The objective sits inside three nested claim stacks. Each stack adds premises
and supports a different conclusion.

\begin{table}[t]
\centering
\small
\begin{tabular}{@{}p{0.24\linewidth}p{0.33\linewidth}p{0.34\linewidth}@{}}
\toprule
Stack & Added premise & Manifesto-facing conclusion \\
\midrule
Preservation & C1/C2/C3 hold on the realized C-Tree for \(f^\ast\) &
Node and root states can replace their raw manifesto spans for
oracle-compatible readouts. \\
Optimization & Preferences or losses depend on the input only through the
preserved oracle interface &
Training on valid root summaries has the same population target as training
on full manifestos. \\
Certificate & Approximate local-law distortion is bounded or estimated from
sampled nodes with logged propensities &
Approximate preservation becomes a finite-sample bound on downstream
objective drift. \\
\bottomrule
\end{tabular}
\caption{\textbf{The theorem ladder.} Each tier keeps the Manifesto target
explicit: the oracle is the policy measurement the analyst wants to preserve.}
\label{tab:v8-theorem-ladder}
\end{table}

The preservation stack is Theorem~\ref{thm:v8-root-preservation} and
Corollary~\ref{cor:v8-preferences-pass-to-root}. It is the local-to-global
claim: if every realized state remains in the same \(f^\ast\)-fiber as the
span it represents, then the root state remains in the same fiber as the whole
manifesto.

The optimization stack adds an oracle-compatibility premise for the downstream
training object. For pairwise preference learning, a comparison rule is
compatible when replacing a manifesto by an equivalent root state does not
change the preferred answer. This covers DPO-style pairwise objectives
\citep{BradleyTerry1952,RafailovEtAl2023DPO} and GRPO-style group objectives
only when the prompt, response generator, policy, ranker, or reward channel
depends on the manifesto through the preserved policy oracle. A preference for
formatting, rhetoric, or style outside the policy target needs a wider oracle
interface.

\begin{proposition}[Manifesto Preference Objective Equivalence]
\label{prop:v8-manifesto-preference-equivalence}
Fix a population of manifestos and a realized C-Tree construction for each
manifesto. If C1/C2/C3 hold exactly for \(f^\ast\), and if the downstream
preference loss is oracle-compatible in the input argument, then replacing each
manifesto by its C-Tree root state leaves the population preference objective
unchanged.
\end{proposition}

The proposition says that root summaries can be used as training inputs only
for preference tasks whose input dependence is already captured by the
manifesto policy target. Exact local laws handle the compression step; the
oracle-compatibility premise handles the downstream objective.

\subsection{Approximate Preservation}

Exact local laws are the clean theorem endpoint. The deployed LLM setting is
approximate. Let \(\Delta_R\) denote the realized mean oracle distortion over
the node population used for the certificate. The subscript \(R\) marks the
represented-span comparison: the state is compared with the raw span it
represents. A method-specific transport constant \(C_{\mathrm{meth}}\) turns
that distortion into downstream objective drift:
\begin{equation}
\label{eq:v8-population-gap}
  |G_{\mathrm{meth}}|
  \le
  C_{\mathrm{meth}}\,\Delta_R .
\end{equation}
For DPO, one sufficient route is a policy-Lipschitz condition for the
log-probability ratios. For group-ranking or reinforcement-style objectives,
the corresponding constant packages the ranker, reward, clipping, and KL
terms. The constant belongs to the downstream training method; the C-Tree audit
supplies \(\Delta_R\).

There are two routes to a usable \(\Delta_R\) bound. A representation route
starts from a function-class approximation claim and transfer moduli, then
turns operator error into C1/C2/C3 budgets. The operational route samples
realized manifesto tree nodes and estimates the actual distortion. The next
section gives the operational route, because it is the route that matters for
deployed manifesto measurement. Appendix~\ref{app:v8-full-proofs} gives the
proof details for the preservation, preference-transport, and certificate
statements used in this ladder.
